\documentclass[11pt,a4paper]{article}
\usepackage[utf8]{inputenc}
\usepackage[T1]{fontenc}
\usepackage[portuguese]{babel}
\usepackage{xcolor}
\usepackage{hyperref}
\usepackage{geometry}
\usepackage{tcolorbox}

\geometry{margin=2.5cm}

\definecolor{primary}{RGB}{39, 174, 96}
\definecolor{secondary}{RGB}{44, 62, 80}
\definecolor{codebg}{RGB}{240, 240, 240}

\hypersetup{colorlinks=true, linkcolor=primary, urlcolor=primary}

\title{\color{secondary}\Huge\textbf{Solução: Exercício 4.7 - Passos de bebê para GitOps}}
\author{\color{primary}DevOps com Kubernetes -- 2026}
\date{}

\begin{document}
\maketitle

\section*{\color{primary}Guia de Solução}
\begin{tcolorbox}[colback=green!5!white,colframe=green!60!black,title=Resolução Passo a Passo]
### Passo a Passo

1. \textbf{Instalar o ArgoCD}: Execute os manifestos de instalação do ArgoCD no cluster.
2. \textbf{Organizar o Repositório}: Certifique-se de que os arquivos YAML do \textit{Log output} estão centralizados em uma pasta no seu repositório Git, idealmente estruturados com um \texttt{kustomization.yaml}.
3. \textbf{Criar a Aplicação no ArgoCD}: Através da interface web do ArgoCD ou de um manifesto do tipo \texttt{Application}, crie a entidade apontando para o seu repositório Git e para o diretório dos manifestos.
4. \textbf{Habilitar Sync Automático}: Nas configurações do ArgoCD \texttt{Application}, mude o \textit{Sync Policy} para \texttt{Automated}. Qualquer \textit{commit} realizado na ramificação principal alterará automaticamente o estado dos \textit{Pods} no Kubernetes.
\end{tcolorbox}

\end{document}
